%% LyX 2.4.4 created this file.  For more info, see https://www.lyx.org/.
%% Do not edit unless you really know what you are doing.
\documentclass[english]{article}
\usepackage[T1]{fontenc}
\usepackage[latin9]{inputenc}
\usepackage{units}
\usepackage{enumitem}
\usepackage{amsmath}
\usepackage{amsthm}
\usepackage{cancel}
\usepackage{geometry}
\geometry{verbose,tmargin=1in,bmargin=1in,lmargin=1in,rmargin=1in}

\makeatletter
%%%%%%%%%%%%%%%%%%%%%%%%%%%%%% Textclass specific LaTeX commands.
\numberwithin{equation}{section}
\numberwithin{figure}{section}
\newlength{\lyxlabelwidth}      % auxiliary length 

%%%%%%%%%%%%%%%%%%%%%%%%%%%%%% User specified LaTeX commands.
\usepackage{fancyhdr}
\usepackage{lastpage}
\pagestyle{fancy}
\usepackage{enumitem}
\usepackage{ifthen}
%sets math fonts to be more readable
\everymath=\expandafter{\the\everymath\displaystyle}
%\DeclareMathSizes{10}{10}{10}{10}
\usepackage{multicol}

\usepackage{tikz}
\usetikzlibrary{plotmarks}
\usepackage{pgfplots}
\pgfplotsset{compat=newest}

\usepackage{calc}
\usepackage[nomessages]{fp}

\usepackage{advdate}    % Advancing/saving dates
\usepackage{datetime}   % Dates formatting
\usepackage{datenumber} % Counters for dates
% Added by lyx2lyx
\setlength{\parskip}{\smallskipamount}
\setlength{\parindent}{0pt}

\makeatother

\usepackage{babel}
\begin{document}
\renewcommand{\author}{Dr.\,James Ashe}

\newcommand{\classNum}{131}

\newcommand{\classSec}{A}

\newcommand{\nameType}{Examples}

\newcommand{\examType}{Worksheet 1.5}

\renewcommand{\date}{\formatdate{9}{9}{2013}}

%%First page's header%%%%%%%%%%%%%%%%%%%%%%
\fancypagestyle{firststyle} %%header settings for first pagr
{
	\fancyhead[L]{MTH \classNum{}(\classSec{}) \author{}\\
	\textbf{\examType{}} \date{}}
	\fancyhead[C]{\textbf{\nameType}}
	\setlength{\headsep}{14pt}
}
%%%%%%%%%%%%%%%%%%%%%%%%%%%%%%%%%%%%%%%%%%

%%Next page's header%%%%%%%%%%%%%%%%%%%%%%%
\setlist{nolistsep}
\lhead{\classNum{}(\classSec{})} 
\chead{\textbf{\nameType}} 
\cfoot{\thepage{}~of~\pageref{LastPage}} 
\setlength{\headsep}{5pt} 
%%%%%%%%%%%%%%%%%%%%%%%%%%%%%%%%%%%%%%%%%%%

%%Various settings; see the preamble for other settings%%%%%
%%\setenumerate[0]{leftmargin=12pt,itemindent=0pt,labelwidth=0pt}
\renewcommand{\labelenumi}{\textbf{\arabic{enumi}.}}
\renewcommand{\labelenumii}{\textbf{\alph{enumii}.}}
\thispagestyle{firststyle}
%%%%%%%%%%%%%%%%%%%%%%%%%%%%%%%%%%%%%%%%%%
\newcommand{\xmax}{10}
\newcommand{\xmin}{-10}
\newcommand{\xstep}{0} %must be xmin+n
\newcommand{\ymax}{10}
\newcommand{\ymin}{-10}
\newcommand{\ystep}{0} %must be ymin+n

\newcommand{\Dspace}{\vspace{1cm}}\textbf{You must show your work to receive credit.}
\begin{enumerate}[resume]
\item Solve the equation using the the square root property. $(5x)^{2}-10=0$.

\begin{enumerate}[resume]
\item []
\begin{eqnarray*}
(5x)^{2} & = & 10\\
\sqrt{(5x)^{2}} & = & \pm\sqrt{10}\\
5x & = & \pm\sqrt{10}\\
x & = & \nicefrac{\pm\sqrt{10}}{5}
\end{eqnarray*}
\vfill
\end{enumerate}
\item Solve the equation by factoring. $2x^{2}-6x-20=0$.

\begin{enumerate}[resume]
\item []
\begin{eqnarray*}
2(x^{2}-3x-10) & = & 0\\
2(x-5)(x+2) & = & 0\\
 &  & x=5\mbox{ and }x=-2
\end{eqnarray*}
\vfill
\end{enumerate}
\item Solve the equation by completing the square. $x^{2}+4x+1=0$.

\begin{enumerate}[resume]
\item []$b=4$ so $\left(\frac{b}{2}\right)^{2}=\left(\frac{4}{2}\right)^{2}=4$
\begin{eqnarray*}
x^{2}+4x & = & -1\\
x^{2}+4x+4 & = & -1+4\\
(x+2)^{2} & = & 3\\
\sqrt{(x+2)^{2}} & = & \pm\sqrt{3}\\
x+2 & = & \pm\sqrt{3}\\
x & = & -2\pm\sqrt{3}
\end{eqnarray*}
\vfill
\end{enumerate}
\item Solve the equation using the quadratic formula. $3x^{2}-3x=4$.

\begin{enumerate}[resume]
\item []First write the equation as, $3x^{2}-3x-4=0$. So$a=3$, $b=-3$,
and $c=-4$.
\begin{eqnarray*}
x & = & \frac{-(-3)\pm\sqrt{(-3)^{2}-4(3)(-4)}}{2(3)}\\
 & = & \frac{3\pm\sqrt{57}}{6}\mbox{ (NOT }\frac{1\pm\sqrt{57}}{2}\mbox{)}
\end{eqnarray*}
\vfill
\end{enumerate}
\end{enumerate}

\end{document}
